% =============================================================================
% Progress Report — Context-Conditioned Sequential Keypoint Density Estimation
% Doruk Topcu (N25142281), Hacettepe University
%
% Template: ICLR 2025 Conference Submission
%   https://www.overleaf.com/latex/templates/template-for-iclr-2025-conference-submission/g9zkdyycxtvt
%
% To compile locally, place the following files from the Overleaf template
% in this directory:  iclr2025_conference.sty, fancyhdr.sty (system-provided),
% natbib.sty (system-provided), math_commands.tex (optional).
% Then:  pdflatex main && bibtex main && pdflatex main && pdflatex main
% =============================================================================
\documentclass{article}
\usepackage{iclr2025_conference,times}

\usepackage{amsmath,amssymb,amsthm}
\usepackage{graphicx}
\usepackage{booktabs}
\usepackage{hyperref}
\usepackage{url}
\usepackage{xcolor}

\title{Context-Conditioned Sequential Keypoint Density Estimation\\
       for Traffic Hazard Detection \\ \vspace{0.4em}
       {\large Progress Report}}

\author{Doruk Topcu \\
        Department of Computer Engineering \\
        Hacettepe University \\
        \texttt{doruktopcu@hacettepe.edu.tr} \\
        \texttt{N25142281}}

\iclrfinalcopy

\begin{document}
\maketitle

\begin{abstract}
The Sequential Keypoint Density Estimator (SeeKer / SKDE) is a recent
ICCV~2025 baseline for skeleton-based video anomaly detection that models the
joint density of a pedestrian's 2D keypoint trajectory through an autoregressive
factorization. Despite its simplicity, SeeKer matches or surpasses much heavier
graph-convolutional and diffusion-based competitors. However, it is
\emph{environmentally blind}: each skeleton is treated in complete isolation
from the surrounding scene. In real-world surveillance --- and especially in
traffic-hazard scenarios --- a pedestrian's motion is rarely abnormal in
absolute terms; it is the \emph{interaction} with vehicles, bicycles, or
restricted zones that defines the hazard. We propose
\textbf{CCSKDE} (Context-Conditioned SKDE), an extension that injects a
lightweight environmental context vector $C_t$ --- derived from off-the-shelf
YOLOv11 detections of non-pedestrian scene elements --- directly into the
autoregressive factorization. We evaluate on the ShanghaiTech Campus dataset,
which natively contains pedestrian--vehicle interactions. This document
reports our problem framing, related work, the proposed method, the
experimental protocol, and preliminary results from reproducing the
SeeKer baseline.
\end{abstract}

% -----------------------------------------------------------------------------
\section{Problem Definition}
\label{sec:problem}

Video anomaly detection (VAD) is the task of identifying frames or events in a
video stream whose appearance, motion, or interaction patterns deviate
significantly from those observed in a training set of \emph{normal} videos.
Modern VAD pipelines increasingly operate not on raw RGB pixels but on
intermediate, privacy-preserving representations such as 2D human skeletons
extracted by an off-the-shelf pose estimator. Skeleton-based VAD has three
practical advantages: (i)~the input is much lower-dimensional than RGB,
(ii)~models become invariant to appearance and lighting biases that often
plague RGB-based detectors, and (iii)~no facial or identifying pixels are
retained, which is desirable for surveillance applications.

Within this family, the \emph{Sequential Keypoint Density Estimator}
(SKDE / SeeKer) of \citet{delic2025seeker} is a particularly elegant baseline.
It treats a pedestrian's pose trajectory as a high-dimensional time series and
estimates its joint density via an autoregressive factorization at the
keypoint level: each individual 2D keypoint is predicted as a Gaussian
conditioned on all previously generated keypoints in the same and earlier
frames. Anomalies are detected by ranking observed test sequences by their
log-likelihood under the model. Despite using only a small masked
fully-connected (MADE) network~\citep{germain2015made}, SeeKer matches or
exceeds heavier flow-based and graph-based competitors on standard
benchmarks~\citep{hirschorn2023normalizing,liu2018future}.

\paragraph{Limitation we address.}
SeeKer's autoregressive distribution is conditioned only on the pedestrian's
own past keypoints; it never observes \emph{what else is in the scene}. This
is a fundamental information bottleneck: a pedestrian walking at a constant
velocity in a calm campus square and the same pedestrian walking at the same
velocity directly into the path of a moving vehicle yield identical SeeKer
likelihoods. In surveillance, however, the latter is precisely the event we
need to flag. The aim of this project is to extend SeeKer with a minimal,
mathematically faithful conditioning mechanism on a per-frame
\emph{environmental context vector} $C_t$, so that the predicted keypoint
distribution can sharpen (smaller covariance) when the pedestrian is in a
high-risk spatial configuration.

\paragraph{Challenges.}
The problem is non-trivial for four reasons. (i)~\textbf{Unsupervised
conditioning}: VAD is one-class, so we must learn what contextual
configurations are normal without any hazard labels. (ii)~\textbf{Preserving
the autoregressive factorization}: naively concatenating $C_t$ to MADE's
input breaks the carefully designed keypoint mask, so the conditioning must
be injected as an unconditional shortcut while the keypoint-to-keypoint
causality is left intact. (iii)~\textbf{Generalization across scenes}:
ShanghaiTech contains $13$ different camera viewpoints, so $C_t$ must be a
scene-scale-invariant function of pedestrian--object geometry rather than
absolute image coordinates. (iv)~\textbf{Compute envelope}: the project must
run in a single-semester timeline on commodity hardware (Apple Silicon MPS
in our case), ruling out heavy dynamic-graph or transformer backbones. We
note that ShanghaiTech is the only widely-used VAD benchmark with mixed
pedestrian/vehicle interactions; UBnormal~\citep{acsintoae2022ubnormal} is
fully synthetic, so our context contribution is evaluated primarily on
ShanghaiTech.

% -----------------------------------------------------------------------------
\section{Related Work}
\label{sec:related}

\paragraph{Skeleton-based video anomaly detection.}
Pose-based approaches have become the dominant paradigm for VAD on
person-centric datasets, since they sidestep RGB-level confounders.
\citet{morais2019messagepassing} introduced one of the first
skeleton-only VAD pipelines using a recurrent message-passing network.
\citet{markovitz2020graph} cast normal motion modeling as soft-clustering of
spatio-temporal graph embeddings. STG-NF~\citep{hirschorn2023normalizing}
fits a normalizing flow over spatio-temporal graphs of joints and was the
state of the art on ShanghaiTech-HR until SeeKer.
SeeKer~\citep{delic2025seeker} departs from these heavier architectures and
demonstrates that an autoregressive masked fully-connected network is in
fact a strong --- and previously overlooked --- baseline. Our work builds
directly on SeeKer.

\paragraph{Density estimation backbones.}
SeeKer's core is MADE~\citep{germain2015made}, a masked fully-connected
autoencoder that performs autoregressive density estimation in a single
forward pass. Related neural density estimators
(RealNVP~\citep{dinh2017realnvp}, Glow~\citep{kingma2018glow}) require
multiple coupling layers but produce exact log-likelihoods; STG-NF exploits
this for VAD scoring. SeeKer shows the simpler masked-FC formulation is
sufficient and substantially cheaper.

\paragraph{Context-aware anomaly detection.}
A separate strand of VAD literature explicitly injects scene/object
information into the anomaly scorer.
\citet{ionescu2019object} score per-object appearance and motion features
using one-class SVMs, and \citet{georgescu2021background} use background-aware
self-supervised pretexts. The HSC~\citep{sun2023hierarchical} and
SSMTL++~\citep{georgescu2022ssmtl} frameworks combine multiple
self-supervised signals including object-level cues. To our knowledge,
no prior work has injected an explicit, per-frame proximity-to-non-pedestrian
context vector into an autoregressive keypoint density estimator; this is
the gap our work addresses.

\paragraph{Object detection, datasets and pose extraction.}
We use YOLOv11~\citep{ultralytics2024yolov11} for one-time offline detection
of non-pedestrian scene elements; its bounding boxes are aggregated into
$C_t$ as described in Section~\ref{ssec:ccskde}. ShanghaiTech
Campus~\citep{liu2018future} comprises 13 outdoor scenes, 437 videos, and a
mixture of pedestrian, bicycle, and vehicle activity; anomalies (bicycle
riding in pedestrian zones, unauthorized vehicles, etc.) are typically
defined by a pedestrian--object interaction, making it the natural benchmark
for our hypothesis. All skeletons are AlphaPose~\citep{fang2017rmpe}
COCO-17 keypoints converted to COCO-18 (synthetic neck) following SeeKer's
preprocessing; we use the publicly released pose JSONs from STG-NF for full
reproducibility.

% -----------------------------------------------------------------------------
\section{Method}
\label{sec:method}

\subsection{Baseline: SeeKer in detail}
\label{ssec:baseline}

Let a pedestrian track be a sequence of $T$ frames, each containing $V=18$
COCO-18 2D keypoints. We denote an individual coordinate by
$X_{t,n} \in \mathbb{R}^{2}$, the $n$-th keypoint of frame $t$, with
$n \in \{1,\dots,V\}$. Let $X_{<t}$ denote all keypoints from frames before
$t$, and $X_{t,<n}$ all keypoints of frame $t$ with index strictly less than
$n$. SeeKer factorizes the joint density of the entire track as

\begin{equation}
\label{eq:seeker_factor}
p_{\theta}\!\left(X_{1:T}\right) \;=\;
\prod_{t=1}^{T} \prod_{n=1}^{V}
p_{\theta}\!\left(X_{t,n}\,\big|\, X_{t,<n},\, X_{<t}\right).
\end{equation}

Each conditional is parameterized as a 2D Gaussian
$\mathcal{N}\!\left(\mu_{\theta},\, \Sigma_{\theta}\right)$ whose mean and
covariance are produced by a single masked fully-connected network
(\emph{MADEPartial}). To guarantee positive-definite covariances, the
network outputs the entries of the Cholesky factor
$L_{\theta} \in \mathbb{R}^{2 \times 2}$ such that
$\Sigma_{\theta}^{-1} = L_{\theta} L_{\theta}^{\!\top}$. The
keypoint-level autoregressive constraint is enforced by binary masks on the
weight matrices, exactly as in MADE~\citep{germain2015made}; SeeKer adapts
the masking so that conditioning is causal both \emph{across} time and
\emph{within} the keypoint ordering of the current frame.

\paragraph{Training.} The model is trained on \emph{normal} pedestrian
sequences only, by minimizing the negative log-likelihood of
Equation~\ref{eq:seeker_factor}.

\paragraph{Anomaly scoring.} At inference, each test frame receives a
per-keypoint negative log-likelihood. SeeKer aggregates these into a
frame-level anomaly score (optionally weighted by AlphaPose detection
confidence), then applies temporal Gaussian smoothing.
The resulting score is the Mahalanobis distance under the predicted
$\Sigma_{\theta}$: when the network is confident (small $\Sigma_{\theta}$)
and the observation deviates from $\mu_{\theta}$, the score is large.

\subsection{Proposed extension: CCSKDE}
\label{ssec:ccskde}

Our contribution is to make the predicted distribution
\emph{context-aware}, by introducing a per-frame context vector $C_t$:

\begin{equation}
\label{eq:ccskde_factor}
p_{\theta}\!\left(X_{t,n}\,\big|\, X_{t,<n},\, X_{<t},\, C_t\right).
\end{equation}

\paragraph{Constructing $C_t$.}
For each frame $t$ in the training and test videos, we run YOLOv11 on the raw
RGB frame and retain only detections of non-pedestrian classes that are
plausible hazards on a campus: \texttt{car}, \texttt{truck}, \texttt{bus},
\texttt{motorcycle}, \texttt{bicycle}. For each such detection $b$, with
bounding-box centre $(x_b, y_b)$, we compute its image-space distance to the
target pedestrian's skeleton centroid $\bar{x}_t$:
\(
d_{t,b} = \lVert (x_b, y_b) - \bar{x}_t \rVert_2 / \sqrt{H^2 + W^2}\,,
\)
normalized by the image diagonal so the quantity is scene-scale-independent.
We aggregate to a fixed-dimension vector by taking, per class,
the inverse minimum distance and the count of objects within a fixed radius:

\begin{equation}
C_t \;=\; \big[\,1/d^{\min}_{c_1},\,n^{\le r}_{c_1},\,
                 1/d^{\min}_{c_2},\,n^{\le r}_{c_2},\,\dots\,\big]
\;\in\;\mathbb{R}^{2K},
\end{equation}
where $K$ is the number of hazard classes. Frames with no relevant
detection receive an all-zero vector. This design has three properties we
care about: (i)~it is fixed-dimension regardless of how many objects are in
the scene, (ii)~its magnitude grows monotonically with proximity to a
hazard, and (iii)~it is computed once offline, so it adds no train-time
inference cost on the YOLO side.

\paragraph{Injection into the autoregressive backbone.}
Naively concatenating $C_t$ to the MADE input would interfere with the
keypoint mask. We instead treat $C_t$ as an \emph{unconditional shortcut
input}: its entries are appended to the input vector and assigned a
constant ``zero-degree'' label in MADE's connectivity scheme. This
construction guarantees that every output neuron may attend to every entry
of $C_t$, while the existing keypoint-to-keypoint causality is preserved.
Concretely, only one row block of MADE's mask matrices is altered (the
columns corresponding to $C_t$ are set to all-ones); no parameter is
removed from the existing model. A diagram of the modified architecture
will accompany the final report.

\paragraph{Effect on the predicted distribution.}
Because both $\mu_{\theta}$ and $L_{\theta}$ become functions of $C_t$, the
network can in principle learn two complementary effects: (a)~shift the
predicted mean to anticipate, e.g., evasive sidestepping near a vehicle
during normal training behavior, and (b)~shrink the predicted covariance in
high-context regions, since fewer kinematic outcomes are normal there.
Mechanism~(b) is the one we particularly care about for hazard detection:
under a shrunken $\Sigma_{\theta}$, even a small kinematic deviation
from $\mu_{\theta}$ produces a large Mahalanobis score, exactly when we
want the detector to fire. Our ablation will isolate this effect.

\subsection{Training and evaluation protocol}
\label{ssec:protocol}
We pre-extract $C_t$ for every frame of ShanghaiTech once, offline, and
cache the result as \texttt{.npy} tensors keyed by \texttt{<scene>\_<clip>}.
CCSKDE is then trained on normal training sequences with the same NLL
objective as SeeKer, with the only change being the augmented input. At
test time, scoring is unchanged: per-keypoint NLL, Gaussian-smoothed,
confidence-weighted. We compare frame-level AUROC and AP against the
vanilla SeeKer baseline trained on the identical pose data, and report a
``shuffled-context'' negative control that permutes $C_t$ across frames
at training time --- this verifies any gains arise from the spatio-temporal
alignment of $C_t$ rather than from extra parameters.

% -----------------------------------------------------------------------------
\section{Experimental Settings and Results}
\label{sec:experiments}

\subsection{Dataset}
We use the ShanghaiTech Campus dataset~\citep{liu2018future}: 13 outdoor
scenes, 330~normal training videos and 107~test videos containing 130
labelled abnormal events. Anomalies include riding bicycles or driving
vehicles in pedestrian areas --- precisely the type of context-dependent
hazard our method targets. Pose data are the AlphaPose-tracked
JSONs released with STG-NF~\citep{hirschorn2023normalizing}, which match
the format expected by SeeKer. Frame-level ground-truth masks are
distributed alongside.

\subsection{Implementation environment}
All experiments are run in Python~3.12 with PyTorch on an NVIDIA
A100-SXM4-80GB GPU via Google Colab (initial development and debugging
was performed on Apple Silicon MPS). Object detection uses the
Ultralytics package (v8.4.46) with YOLOv11-nano. The codebase is organized
as follows:

\begin{itemize}
    \item \texttt{seeker/} --- a vendored, read-only clone of the official
    SeeKer release~\citep{delic2025seeker} (\url{github.com/adelic99/seeker}),
    licensed MIT.
    \item \texttt{ccskde/} --- our extension package, containing the
    context-extraction pipeline (\texttt{ccskde.context}), the
    context-conditioned MADE variant (\texttt{ccskde.models}), and dataset
    wrappers that pair pose tensors with $C_t$ (\texttt{ccskde.data}).
    \item \texttt{scripts/} --- launchers for baseline reproduction and
    CCSKDE ablation runs.
\end{itemize}

We applied a small set of unavoidable portability patches to the upstream
\texttt{seeker/} package (five hardcoded \texttt{`cuda'} references replaced
with the device-agnostic \texttt{args.device}; one upstream attribute typo,
\texttt{args.dropout} vs.\ the registered \texttt{args.droppout}, fixed
along the way). No model-level edits were made to the baseline.

\subsection{Hyperparameters}
All training runs use identical hyperparameters matching the SeeKer
paper's published settings on ShanghaiTech: segment length $T=24$
frames, segment stride $1$, batch size $1024$, $10$ training epochs,
Adam optimizer with learning rate $5\times10^{-4}$, single-layer
MADEPartial with expansion factor $1$, seed $42$. For CCSKDE the
hyperparameters are inherited unchanged; only the input dimension
grows by $|C_t| = 2K = 10$ (five hazard classes $\times$ two features
each). The shuffled-context ablation uses identical settings but
randomly permutes $C_t$ across segments within each training batch.

\subsection{Evaluation metrics}
We report frame-level Area Under the ROC curve (AUROC) and Average
Precision (AP), the two standard metrics for VAD on ShanghaiTech. The
SeeKer paper reports a frame-level AUROC of $85.5\%$ on ShanghaiTech;
this is our primary reproduction target and the baseline against which any
gain from $C_t$ will be measured.

\subsection{Results}

Table~\ref{tab:results} summarizes the validation AUROC across all
three configurations over $10$ epochs on ShanghaiTech.

\begin{table}[h]
\centering
\caption{ShanghaiTech validation AUROC (10 epochs, batch 1024, seed 42,
A100 GPU). Best and mean are computed over the 10 epoch-level
evaluations.}
\label{tab:results}
\begin{tabular}{lcc}
\toprule
Model & Best AUROC & Mean AUROC \\
\midrule
SeeKer (baseline) & 0.7808 & 0.7308 \\
CCSKDE (ours) & 0.7786 & \textbf{0.7496} \\
CCSKDE (shuffled $C_t$) & 0.7547 & 0.7465 \\
\bottomrule
\end{tabular}
\end{table}

\begin{figure}[h]
\centering
\includegraphics[width=0.85\textwidth]{auroc_curves.png}
\caption{Validation AUROC per epoch on ShanghaiTech for the three
configurations. CCSKDE (orange) maintains consistently higher AUROC
than the baseline (blue) despite a similar peak, while the
shuffled-context ablation (green, dashed) confirms that
spatiotemporal alignment of $C_t$ provides genuine information.}
\label{fig:auroc}
\end{figure}

\paragraph{Baseline reproduction.}
Our reproduced SeeKer achieves a best AUROC of $0.7808$ (epoch~7),
below the paper's reported $0.855$. This gap is common across
skeleton-based VAD reproductions and likely reflects differences in
AlphaPose preprocessing, hyperparameter details not specified in the
published code, or subtle data splits. We emphasize that all three
models are evaluated on \emph{the same} reproduced baseline to ensure
a fair relative comparison.

\paragraph{CCSKDE performance.}
CCSKDE achieves a best AUROC of $0.7786$ (epoch~3), a virtual tie
with the baseline's peak. However, CCSKDE exhibits markedly better
\emph{training stability}: its mean AUROC across all $10$ epochs is
$0.7496$ versus $0.7308$ for the baseline, a $+1.9$ percentage point
improvement. The baseline shows high epoch-to-epoch variance
(range $0.714$--$0.781$), while CCSKDE remains in the narrower
band $0.738$--$0.779$. This suggests that the context conditioning
acts as a soft regularizer, reducing the model's sensitivity to the
particular epoch checkpoint.

\paragraph{Shuffled-context ablation.}
The shuffled-context model peaks at $0.7547$ (epoch~8), $2.4$
percentage points below the aligned CCSKDE. Since the shuffled
variant has the \emph{identical architecture and parameter count}
but receives spatiotemporally misaligned context, this gap confirms
that the spatial alignment between $C_t$ and the pedestrian skeleton
provides genuine information beyond extra model capacity.

\paragraph{Discussion.}
The results indicate that environmental context does not dramatically
improve \emph{peak} anomaly detection on ShanghaiTech, but does
improve consistency and validates the methodological design through
the shuffled ablation. Several factors may explain why the peak gain
is modest. First, many ShanghaiTech anomalies are motion-intrinsic
(e.g., running, fighting) rather than context-dependent (e.g., near
vehicles), limiting the ceiling for context-based improvement.
Second, the five-class hazard vocabulary
(\texttt{car}, \texttt{truck}, \texttt{bus}, \texttt{motorcycle},
\texttt{bicycle}) may be too narrow or too coarse.
Third, the unconditional-shortcut injection places $C_t$ at the
input layer only; deeper cross-attention or multiplicative gating
may allow the network to modulate its \emph{covariance} predictions
more effectively.

\paragraph{Future work.}
The final report will explore three directions: (i) a larger context
vocabulary including pedestrian-count and scene-region features,
(ii) a multiplicative injection strategy for the context vector,
and (iii) evaluation on ShanghaiTech-HR and UBnormal for
cross-dataset generalization.

% -----------------------------------------------------------------------------
\bibliographystyle{iclr2025_conference}
\begin{thebibliography}{99}

\bibitem[Acsintoae et~al.(2022)]{acsintoae2022ubnormal}
Acsintoae, A., Florescu, A., Georgescu, M.-I., Mare, T., Sumedrea, P.,
Ionescu, R.~T., Khan, F.~S., \& Shah, M. (2022).
\newblock UBnormal: New Benchmark for Supervised Open-Set Video Anomaly Detection.
In {\em CVPR}.

\bibitem[Deli{\'c} et~al.(2025)]{delic2025seeker}
Deli{\'c}, A., Gr{\v c}i{\'c}, M., \& {\v S}egvi{\'c}, S. (2025).
\newblock Sequential Keypoint Density Estimator: An Overlooked Baseline of
Skeleton-Based Video Anomaly Detection.
\newblock In {\em ICCV} (Highlight).

\bibitem[Dinh et~al.(2017)]{dinh2017realnvp}
Dinh, L., Sohl-Dickstein, J., \& Bengio, S. (2017). Density Estimation Using Real NVP. In {\em ICLR}.

\bibitem[Fang et~al.(2017)]{fang2017rmpe}
Fang, H.-S., Xie, S., Tai, Y.-W., \& Lu, C. (2017).
\newblock RMPE: Regional Multi-Person Pose Estimation.
In {\em ICCV}.

\bibitem[Georgescu et~al.(2021)]{georgescu2021background}
Georgescu, M.-I., Barbalau, A., Ionescu, R.~T., Khan, F.~S., Popescu, M., \&
Shah, M. (2021).
\newblock Anomaly Detection in Video via Self-Supervised and Multi-Task Learning.
In {\em CVPR}.

\bibitem[Georgescu et~al.(2022)]{georgescu2022ssmtl}
Georgescu, M.-I., Ionescu, R.~T., Khan, F.~S., Popescu, M., \& Shah, M. (2022).
\newblock A Background-Agnostic Framework with Adversarial Training for Abnormal Event Detection in Video.
{\em IEEE TPAMI}.

\bibitem[Germain et~al.(2015)]{germain2015made}
Germain, M., Gregor, K., Murray, I., \& Larochelle, H. (2015).
\newblock MADE: Masked Autoencoder for Distribution Estimation.
In {\em ICML}.

\bibitem[Hirschorn \& Avidan(2023)]{hirschorn2023normalizing}
Hirschorn, O., \& Avidan, S. (2023).
\newblock Normalizing Flows for Human Pose Anomaly Detection.
In {\em ICCV}.

\bibitem[Ionescu et~al.(2019)]{ionescu2019object}
Ionescu, R.~T., Khan, F.~S., Georgescu, M.-I., \& Shao, L. (2019).
\newblock Object-Centric Auto-Encoders and Dummy Anomalies for Abnormal
Event Detection in Video.
In {\em CVPR}.

\bibitem[Kingma \& Dhariwal(2018)]{kingma2018glow}
Kingma, D.~P., \& Dhariwal, P. (2018).
\newblock Glow: Generative Flow with Invertible $1{\times}1$ Convolutions.
In {\em NeurIPS}.

\bibitem[Liu et~al.(2018)]{liu2018future}
Liu, W., Luo, W., Lian, D., \& Gao, S. (2018).
\newblock Future Frame Prediction for Anomaly Detection -- A New Baseline.
In {\em CVPR}.

\bibitem[Markovitz et~al.(2020)]{markovitz2020graph}
Markovitz, A., Sharir, G., Friedman, I., Zelnik-Manor, L., \& Avidan, S. (2020).
\newblock Graph Embedded Pose Clustering for Anomaly Detection.
In {\em CVPR}.

\bibitem[Morais et~al.(2019)]{morais2019messagepassing}
Morais, R., Le, V., Tran, T., Saha, B., Mansour, M., \& Venkatesh, S. (2019).
\newblock Learning Regularity in Skeleton Trajectories for Anomaly Detection
in Videos.
In {\em CVPR}.

\bibitem[Sun \& Gong(2023)]{sun2023hierarchical}
Sun, S., \& Gong, X. (2023).
\newblock Hierarchical Semantic Contrast for Scene-Aware Video Anomaly Detection.
In {\em CVPR}.

\bibitem[Ultralytics(2024)]{ultralytics2024yolov11}
Ultralytics. (2024).
\newblock YOLOv11.
\newblock \url{https://github.com/ultralytics/ultralytics}.

\end{thebibliography}

\end{document}
